\section{Conclusion}

We have presented \textit{Query-GSAS}, a retrieval-augmented generation system
that provides natural-language, citation-linked access to the complete
GSAS-II documentation corpus using only local computation.
The system indexes 5 sets of documentation sources totalling 331 HTML files and $\sim$500,000 words into $\sim$5{,}838 embedded chunks in a persistent ChromaDB vector database.
The system achieves Recall@6 of 0.98 on a representative 40-question
evaluation set, and generates fluent, inline-cited answers in
approximately 8\,s on modern CPU hardware.

The fully local architecture --- built on \texttt{BAAI/bge-base-en-v1.5}
sentence embeddings served via the \texttt{fastembed} ONNX runtime (PyTorch
not required), persistent ChromaDB storage, and either Ollama-served
quantised language models or in-process \texttt{llama-cpp-python}
inference --- satisfies potential data-residency and security requirements
without sacrificing answer quality.
A post-processing citation injector and first-appearance renumbering
scheme ensure that generated answers carry dense, sequentially numbered
inline citations rendered as clickable hyperlinks with a
journal-style bibliography in the web UI, regardless of the underlying
model's instruction-following capability.

Three complementary interfaces (CLI, wxPython desktop dialog, and
web application) allow \textit{Query-GSAS} to serve users from the command line,
from within the GSAS-II graphical interface via a single-function Help menu
integration (\texttt{show\_assistant(parent=self)}), and from a shared
web server accessible to multiple beamline users simultaneously.
The Ollama lifecycle is managed transparently by the GUI dialog, with
\texttt{atexit} fallback handling for the case where GSAS-II terminates
without explicitly closing the assistant window.

The open-source release (\texttt{pip install gsas-query};
\url{https://github.com/AdvancedPhotonSource/Query-GSAS}) is intended
to serve as a working template for researchers and developers who wish
to build similar RAG-based documentation assistants for other large
scientific software packages.
